\documentclass[12pt]{article}
\usepackage{resource/shortex}
\usepackage{mathtools,amssymb}
\usepackage{enumitem}
\usepackage{datetime}
\usepackage{fancyhdr}
\usepackage{ifthen}
\usepackage{pifont}
\usepackage{mathrsfs}
\usepackage[normalem]{ulem}
%\RequirePackage[usenames,dvipsnames]{color}


%% showing solutions 
\newboolean{showsols}
\setboolean{showsols}{false}
\newcommand{\showsolutions}{\setboolean{showsols}{true}}
\newlength{\solspace}
\setlength{\solspace}{16em}
\newcommand{\nosolspace}{\setlength{\solspace}{0em}}

\newcommand{\solution}[2]{\ifthenelse{\boolean{showsols}}{{\par\textcolor{red}{Rubric criterion: #1}\newline\textcolor{blue}{#2}}}{\vspace{\solspace}}}
\newcommand{\npforprint}{\ifthenelse{\boolean{showsols}}{}{\newpage}}
\newcommand{\vsforprint}[1]{\ifthenelse{\boolean{showsols}}{}{\vspace{#1}}}

%% grading
\newcommand{\grades}[1]{~\newline\noindent\framebox[\textwidth][c]{E\quad \quad \quad \quad S \quad \quad \quad \quad R \quad \quad \quad \quad N}}

\newcommand{\questiongrades}{\noindent\emph{Each item is marked $\checkmark$, $\checkmark\!-$, or \ding{55}:}\begin{center}\emph{$\checkmark$ = capable  $~~$ $\checkmark\!-$ = mostly capable  $~~$ \ding{55} = not yet}\end{center} }

\newcommand{\triAssessments}{$\checkmark\quad\checkmark\!-\quad$\ding{55}\quad}
\newcommand{\biAssessments}{$\checkmark\quad$\ding{55}\quad}

%% line for writing name and BUID
\newcommand{\nameBUID}{\noindent\textbf{Name:}\underline{\phantom{XXXXXXXXXXXXXXXXXXX}}\hfill\textbf{BUID:}\underline{\phantom{XXXXXXXXXXXXX}}\newline}


%% formatting 
%\allowdisplaybreaks[3]
\renewcommand{\headrulewidth}{.4mm} % header line width
\renewcommand{\arraystretch}{1.25}


%% page style 
\pagestyle{fancy}
\fancyhf{}
\rhead{Spring 2026}
\lhead{CAS MA 214: Applied Statistics}
%\rfoot{\thepage}

%--------------------
\begin{document}

\begin{center}
\textbf{\Large Project 1 \\ Statistical Analysis and Model Application} \\
\end{center}


\section*{Overview}
This project provides an opportunity to apply the concepts and tools you have learned in this course to a real-world dataset of your choice. Through this experience, you will practice the full process of statistical modeling — from data exploration and model fitting to interpretation and communication of results. The goal is to help you see how regression models can be used to understand relationships, make predictions, and support data-driven decisions in practical contexts. \\ \\
Your group will select a dataset of your interest and conduct modeling using regression models.

\begin{itemize}
    \item Choose a dataset that includes a response variable and one or more predictors.
    \item Model the relationship between the response variable and predictors using regression methods.
    \item A good example would be *Chapter 10.1 Case Study: Houses for Sale*.
\end{itemize}

\textbf{Your group will present the analyses through a pre-recorded video presentation.}  
You will be assigned to watch presentations by other groups through the Blackboard course page.  
Each group member must ask at least three questions about other groups’ presentations and respond to at least three questions about your own.


\section*{Suggested Outline for Presentation}
\begin{itemize}
    \item Introduction
    \begin{enumerate}
        \item Introduction of the dataset and its source
        \item The story of your reasoning
        \item Variables of interest 
    \end{enumerate}
    \item Data Analysis and Modeling
    \begin{enumerate}
        \item Exploratory Data Analysis Part
        \begin{enumerate}
            \item Summary statistics 
            \item The relationship between response variable and predictor(s)  
        \end{enumerate}
        \item Modeling Part
        \begin{enumerate}
            \item Fitting regression model(s)
            \item Choosing the best model (model selection)
            \item Interpretation of the model output(s)
        \end{enumerate} 
        \item Model Diagnostics and Validation
        \begin{enumerate}
            \item Residual analysis
        \end{enumerate}

    \end{enumerate}
    
    \item Conclusion
        \begin{enumerate}
            \item Summarize key insights
            \item Suggestion of future analysis 
        \end{enumerate} 
\end{itemize}

\section*{Project requirements}
\begin{itemize}
    \item Source of data must be cited
    \item R code should be included (Rscript or Rmd files)
    \item References should be provided at the end
\end{itemize}


\section*{Presentation}
\begin{itemize}
    \item Length of the presentation : 5 min 

\end{itemize}



\section*{Possible Data Sources to Consider}
\begin{itemize}
    \item \href{https://www.openintro.org/data/}{OpenIntro Datasets}
    \item \href{https://cran.r-project.org/web/packages/fivethirtyeight/vignettes/fivethirtyeight.html}{fivethirtyeight 
    R Package}
    \item \href{https://stat.ethz.ch/R-manual/R-devel/library/datasets/html/00Index.html}{datasets R package}
    \item \href{https://www.kaggle.com/datasets}{Kaggle datasets}
    \item \href{https://data.gov/}{US government's open data}
\end{itemize}

% \section*{Example Roles by Tasks}
% \begin{enumerate}
%     \item Role 1 : Numerical Data Analysis \\
%     Role 2 : Categorical Data Analysis \\
%     Role 3 : Video Presentation \\
%     Role 4 : Video Editing \\
%     \item Role 1 : Introduction Part   \\
%     Role 2 : Data Analysis Part \\
%     Role 3 : Conclusion Part \\
%     Role 4: Video Editing  \\
%     Each member (Role 1 -3) is responsible for presentating their own part. \\
%     \item Role 1 : Introduction Part   \\
%     Role 2 :  Categorical Data Analysis  Part \\
%     Role 3 : Numerical Data Analysis Part\\
%     Role 4: Conclusion and Video Editing \\
%     Each member is responsible for presentating their own part. \\
% \end{enumerate}

\section*{Deliverables}
\begin{itemize}
    \item Deliverable 1 : Project Outline due on Feb 3
    \item Deliverable 2 : Project Plan (Draft Slides + R code + Dataset) due on Feb 17
    \item Deliverable 3 : Recorded Presentation + Final Presentation Slides due on Mar 3
\end{itemize}

\section*{Important Dates}
\begin{itemize}
    \item Jan 28, Wed : Project 1 logisstics lab 
    \item Feb 3, Tuesday at 10:00 PM : Deliverable 1 - Project Outline due
    \item Feb 11, Wed : Project 1 work lab
    \item Feb 17, Tuesday at 10:00 PM : Deliverable 2 - Project Plan (Draft Slides + R code + Dataset) due
    \item Feb 18, Wed : Project 1 work lab
    \item Feb 25, Wed : Project 1 work lab
    \item Mar 3, Tuesday at 10:00 PM : Deliverable 3 - Recorded Presentation + Final Presentation Slides due
\end{itemize}


\end{document}



